% 60/40 in 2022: When the Stock-Bond Correlation Flipped — research-note PDF.
%
% PROSE AND CODE ARE VERBATIM from the website article:
%   site/src/content/articles/sixty-forty-correlation-flip.tsx
%   site/src/content/articles/sixty-forty-code.ts
% Metadata from lib/articles.ts; project-card fallbacks from lib/topics.ts.
%
% NO CHARTS. The article's header records it as deliberately unexecuted: no
% data module, no computed numbers, prose and code only. The figures quoted in
% the prose are the author's own, written into his markdown.
\documentclass{pyportfolios-note}

\noteshorttitle{60/40 in 2022}

\begin{document}
\thispagestyle{notecover}

\notemasthead{Quantitative Insights}{Case Study}{Q3 2026}{}

\notetitle{60/40 in 2022:\\When the Stock-Bond\\Correlation Flipped}

\notedeck{During the 2022 inflation shock, the balanced portfolio failed because
its diversifier stopped diversifying.}

For two decades, the 60/40 portfolio rested on one quiet assumption:
\textbf{when stocks fall, bonds rally}. In 2022 that assumption broke. Stocks
fell \textasciitilde18\%, long Treasuries fell \textasciitilde31\% — the worst
year for a US balanced portfolio since the Global Financial Crisis, and by some
measures since the 1930s. This case study reconstructs what happened, why the
correlation flipped, and what it means for every portfolio built on the
stock-bond hedge.

\notecard
  {Portfolio Optimization}
  {Pandas · statsmodels · seaborn}
  {SPY + AGG (+ TLT for the long-duration pain)}
  {Jan 2000 – Dec 2023, focus 2022}
  {Stress-testing the diversification assumption every balanced portfolio
   silently relies on}

\begin{notecode}
import numpy as np
import pandas as pd
import matplotlib.pyplot as plt
import seaborn as sns

plt.rcParams["figure.dpi"] = 110
sns.set_style("whitegrid")
\end{notecode}

% ================================ 1 The setting: two decades of a free hedge =
\noteneedroom{170bp}
\notesection{1.\notenumsep The setting: two decades of a free hedge}

\begin{multicols}{2}
From roughly 2000 to 2021, US stock and Treasury returns were \textbf{negatively
correlated}. Every equity selloff — 2008, 2011, 2018, March 2020 — saw bonds
rally as investors fled to safety and the Fed cut rates. A 60/40 investor got
equity upside with a built-in shock absorber, and the strategy compounded
through every crisis.

The subtlety everyone forgot: that negative correlation is a \textbf{regime},
not a law. It held because inflation was low and stable, so growth shocks
dominated — bad news for stocks was good news for bonds (rate cuts coming). In
the 1970s–90s, when \textit{inflation} shocks dominated, the correlation had
been \textbf{positive}.
\end{multicols}

% ================================== 2 The event: 2022 month by month =========
\noteneedroom{150bp}
\notesection{2.\notenumsep The event: 2022 month by month}

\begin{multicols}{2}
Inflation, dismissed as ``transitory'' in 2021, printed 7\%+ into 2022. The Fed
delivered the fastest hiking cycle in four decades — from 0.25\% to 4.50\% in
ten months. Rising yields crushed bond prices (see our duration tutorial: TLT's
\textasciitilde17.5-year duration $\times$ \textasciitilde2\% yield rise
$\approx$ $-$35\%), while the same rate shock compressed equity multiples.

\textbf{Both engines of 60/40 stalled at once.} There was nowhere to hide inside
the classic mix.
\end{multicols}

% ====================== 3 The mechanism: why inflation flips the sign ========
\noteneedroom{170bp}
\notesection{3.\notenumsep The mechanism: why inflation flips the sign}

\begin{multicols}{2}
\begin{notebullets}
\item \textbf{Growth-shock regime} (2000–2021): bad economy $\rightarrow$ stocks
      fall, Fed cuts $\rightarrow$ bonds rise. Correlation \textbf{negative}.
      60/40 self-hedges.
\item \textbf{Inflation-shock regime} (1970s, 2022): inflation up $\rightarrow$
      Fed hikes $\rightarrow$ \textit{discount rates rise for everything}
      $\rightarrow$ stocks \textbf{and} bonds fall together. Correlation
      \textbf{positive}. The hedge becomes a second exposure to the same risk.
\end{notebullets}

One variable — which type of shock dominates — determines whether bonds protect
you or double your bet.
\end{multicols}

% ========================================================= 4 The evidence ====
\noteneedroom{150bp}
\notesection{4.\notenumsep The evidence}

\notesubsection{4.1\notenumsep The data}

SPY (S\&P 500), AGG (aggregate bonds) and TLT (long Treasuries) from AGG's
inception through 2023.

\begin{notecode}
TICKERS = ["SPY", "AGG", "TLT"]
START, END = "2003-10-01", "2023-12-31"

def load_prices(tickers, start, end):
    """Adjusted-close prices: yfinance first, Stooq as fallback."""
    try:
        import yfinance as yf
        df = yf.download(tickers, start=start, end=end, auto_adjust=True, progress=False)["Close"]
        if not df.empty:
            return df[tickers].dropna()
    except Exception as exc:
        print(f"yfinance failed ({exc}); trying Stooq…")
    cols = {}
    for t in tickers:
        url = f"https://stooq.com/q/d/l/?s={t.lower()}.us&i=d"
        cols[t] = pd.read_csv(url, parse_dates=["Date"], index_col="Date")["Close"].rename(t)
    return pd.concat(cols, axis=1).loc[start:end].dropna()

px = load_prices(TICKERS, START, END)
rets = px.pct_change().dropna()
print(f"{len(px)} trading days, {px.index[0].date()} → {px.index[-1].date()}")
\end{notecode}

\notesubsection{4.2\notenumsep Annual returns: 2022 in context}

The one chart that tells the story — find another year where \textit{both} bars
are deeply negative.

\begin{notecode}
annual = (1 + rets).resample("YE").prod() - 1
annual.index = annual.index.year

ax = annual[["SPY", "AGG"]].plot(kind="bar", figsize=(11, 4.5),
                                  color=["steelblue", "darkorange"], width=0.8)
ax.axhline(0, color="black", lw=0.8)
ax.set_ylabel("Total return"); ax.set_title("Annual returns: stocks (SPY) vs bonds (AGG)")
ax.legend(["SPY (stocks)", "AGG (bonds)"])
for yr, row in annual.iterrows():
    if yr == 2022:
        ax.axvspan(list(annual.index).index(yr) - 0.5, list(annual.index).index(yr) + 0.5,
                   color="crimson", alpha=0.10)
plt.tight_layout(); plt.show()

print(annual.loc[[2008, 2020, 2022]].round(3))
print("\n2008 & 2020: bonds cushioned the crash. 2022: they amplified it.")
\end{notecode}

\notesubsection{4.3\notenumsep The correlation flip}

Rolling 1-year correlation of daily stock and bond returns. Two decades below
zero — then 2022.

\begin{notecode}
roll_corr = rets["SPY"].rolling(252).corr(rets["AGG"])

fig, ax = plt.subplots(figsize=(11, 4.5))
ax.plot(roll_corr, color="steelblue", lw=1.2)
ax.axhline(0, color="black", lw=0.8)
ax.axvspan(pd.Timestamp("2022-01-01"), pd.Timestamp("2023-01-01"), color="crimson", alpha=0.12)
ax.set_title("Rolling 1-year stock-bond correlation (SPY vs AGG)")
ax.set_ylabel("Correlation")
plt.tight_layout(); plt.show()

print(f"Average correlation 2004–2021: {roll_corr.loc[:'2021'].mean():+.2f}")
print(f"Peak correlation in 2022–23:   {roll_corr.loc['2022':].max():+.2f}")
\end{notecode}

\notesubsection{4.4\notenumsep The damage: 60/40 drawdown}

A monthly-rebalanced 60/40 (SPY/AGG). Compare the 2022 drawdown with the GFC —
and note how much \textit{faster} 2022 hurt, because nothing offset anything.

\begin{notecode}
w = {"SPY": 0.60, "AGG": 0.40}
port_rets = (rets[list(w)] * pd.Series(w)).sum(axis=1)          # daily, approx monthly rebalance
port_val = (1 + port_rets).cumprod()
drawdown = port_val / port_val.cummax() - 1

fig, ax = plt.subplots(figsize=(11, 4.5))
ax.fill_between(drawdown.index, drawdown * 100, 0, color="steelblue", alpha=0.6)
ax.axvspan(pd.Timestamp("2022-01-01"), pd.Timestamp("2023-01-01"), color="crimson", alpha=0.12)
ax.set_title("60/40 (SPY/AGG) drawdown")
ax.set_ylabel("Drawdown (%)")
plt.tight_layout(); plt.show()

dd_2022 = drawdown.loc["2022":"2023"].min()
dd_gfc  = drawdown.loc["2007":"2010"].min()
ret_2022 = (1 + port_rets.loc["2022"]).prod() - 1
print(f"60/40 total return 2022:  {ret_2022:+.1%}")
print(f"Max drawdown 2022–23:     {dd_2022:+.1%}")
print(f"Max drawdown GFC 2008-09: {dd_gfc:+.1%}")
\end{notecode}

And the counterfactual that stings: a 60/40 built with \textbf{TLT} instead of
AGG — more duration, more ``hedge'' — did \textit{worse} in 2022, because the
hedge asset itself was the epicenter.

\begin{notecode}
for bond in ["AGG", "TLT"]:
    pr = 0.60 * rets["SPY"] + 0.40 * rets[bond]
    r22 = (1 + pr.loc["2022"]).prod() - 1
    print(f"60/40 with {bond}: 2022 return {r22:+.1%}")
\end{notecode}

% ========================================================= 5 Post-mortem =====
\noteneedroom{330bp}
\notesection{5.\notenumsep Post-mortem}

\begin{multicols}{2}
\notesubsection{What failed}
\begin{notebullets}
\item \textbf{The model, not the math} — MVO and risk parity both treated the
      stock-bond correlation as a stable input. It's regime-dependent, and the
      regime variable is \textit{inflation}.
\item \textbf{Duration as a hedge} — in an inflation shock, duration is the
      \textit{risk}, not the hedge (see the bond tutorial's 2022 ETF table).
\item \textbf{Recency} — twenty years of negative correlation felt like a law of
      nature. The 1970s said otherwise all along.
\end{notebullets}

\notesubsection{What survived}
\begin{notebullets}
\item \textbf{Commodities and trend-following} had a banner 2022 — the
      diversifiers nobody wanted during the long bull market.
\item \textbf{Short-duration bonds} (SHY) lost little — the failure was
      duration, not bonds per se.
\item \textbf{The principle of diversification} — but across \textit{risk
      regimes} (inflation vs growth), not just asset classes.
\end{notebullets}

\notesubsection{The lessons}
\begin{notebulletsnum}
\item Correlation inputs deserve the same stress-testing as returns —
      \textit{conditional} on inflation regime, not unconditional averages.
\item A hedge that depends on a regime is a position on that regime persisting.
\item The fix isn't abandoning 60/40 — it's knowing which environments it
      insures against, and which it doesn't.
\end{notebulletsnum}

\notesubsection{Related content}
\begin{notebullets}
\item \textbf{Bond Pricing, Duration \& Convexity} — why TLT lost 31\% (the
      mechanics behind this case)
\item \textbf{Risk Parity from Scratch} — the strategy this correlation flip
      hurt most
\item \textbf{Copulas \& Tail Dependence} — modelling co-movement beyond a
      single correlation number
\end{notebullets}
\end{multicols}

\end{document}
